\section{Evaluation at the Pell scale and the constitutive Barbero--Immirzi functional}
\label{iii:sec:pell-barbero}

The relationship between the vacuum response and the Catalan moment allows
two levels of transport to be composed: the weighted depth of the
quarter-turn character and evaluation of the angular channels. The Pell scale
distinguishes an argument of the moment collection; the vacuum channels,
by contrast, receive the arguments determined by the angular pair.
The Barbero--Immirzi functional is expressed on the latter through
the same constitutive response and its inverse. The connection is thus established
between defined operations, preserving the position of each evaluation
within the APP--TRIT--TPK construction.

\subsection{Oriented moment and evaluation at the Pell unit}

Write
\begin{equation}
 G_{\rm Cat}:=\mathcal C_2(1)
 =\sum_{k\geq0}\frac{(-1)^k}{(2k+1)^2},\qquad
 \lambda_{\rm P}=2+\sqrt3,\qquad
 \rho_{\rm P}=2-\sqrt3=\lambda_{\rm P}^{-1}.
 \label{iii:eq:pell-scales}
\end{equation}
The constant $G_{\rm Cat}$ is the endpoint value of the moment constructed in
\eqref{iii:eq:c2}. The pair $\lambda_{\rm P},\rho_{\rm P}$ corresponds to
the expanding and contracting orientations of the Pell unit:
$(2+\sqrt3)(2-\sqrt3)=1$. The norm identity $2^2-3\cdot1^2=1$
fixes reciprocity. Specialization of the moment at $\rho_{\rm P}$
follows construction of its oriented character.

Differential equation \eqref{iii:eq:c2-integral}, with zero condition
at the origin, gives
\begin{equation}
 \mathcal D\mathcal C_2(s)=\arctan s,\qquad
 \mathcal C_2(s)=\int_0^s\frac{\arctan t}{t}\,dt,
 \qquad \mathcal D=s\frac{d}{ds}.
 \label{iii:eq:pell-moment-integral}
\end{equation}
Indeed, integrating $\mathcal D^2\mathcal C_2=s/(1+s^2)$ with respect to
$ds/s$ produces $\arctan s$; a second integration produces the stated
expression. The integrand extends continuously to $t=0$ with value $1$.
The angular form of this integral allows exact determination of the
evaluation at $\rho_{\rm P}$.

\begin{lemma}[Angular integral of the odd character]
\label{iii:lem:log-tangent}
For $0<\theta\leq\pi/4$,
\begin{equation}
 -\int_0^\theta\log(\tan x)\,dx
 =\sum_{\substack{n\geq1\\n\text{ odd}}}
      \frac{\sin(2n\theta)}{n^2}.
 \label{iii:eq:log-tangent}
\end{equation}
\end{lemma}
\begin{proof}
For $0<a<1$, the logarithm expansion inside the unit disk implies
\begin{align*}
 L_a(x)
 &:=\log|1-ae^{2ix}|-\log|1+ae^{2ix}|\\
 &=-2\sum_{\substack{n\geq1\\n\text{ odd}}}
       \frac{a^n\cos(2nx)}{n}.
\end{align*}
The series converges uniformly in $x$ for every $a<1$; termwise
integration gives
\[
 -\int_0^\theta L_a(x)\,dx
 =\sum_{\substack{n\geq1\\n\text{ odd}}}
       \frac{a^n\sin(2n\theta)}{n^2}.
\]
As $a\uparrow1$, the quotient of moduli converges to
$|1-e^{2ix}|/|1+e^{2ix}|=\tan x$ for $0<x\leq\theta$.
The logarithm's singularity at the origin is integrable. To justify
the limit, one may restrict to $a\geq1/2$ and observe that
\[
 |1-ae^{2ix}|^2=(1-a)^2+4a\sin^2x
 \geq2\sin^2x,
 \qquad 1\leq|1+ae^{2ix}|\leq2
\]
on $0\leq x\leq\pi/4$. Together with the upper bound
$|1-ae^{2ix}|\leq2$, these inequalities give
$|L_a(x)|\leq c+|\log x|$ with a constant independent of $a$.
Dominated convergence permits passage to the limit in the integral.
For the integrated series, domination by $\sum n^{-2}$ applies.
The two limits prove \eqref{iii:eq:log-tangent}.
\end{proof}

\begin{theorem}[Catalan--Pell evaluation]
\label{iii:thm:catalan-pell}
The oriented moment satisfies
\begin{align}
 \mathcal C_2(\rho_{\rm P})
 &=\frac23G_{\rm Cat}-\frac{\pi}{12}\log\lambda_{\rm P},
 \label{iii:eq:catalan-pell}\\
 \mathcal D\mathcal C_2(\rho_{\rm P})&=\frac{\pi}{12},&
 \mathcal D^2\mathcal C_2(\rho_{\rm P})&=\frac14.
 \label{iii:eq:catalan-pell-derivatives}
\end{align}
\end{theorem}
\begin{proof}
The substitution $t=\tan x$ in \eqref{iii:eq:pell-moment-integral},
followed by integration by parts, gives
\begin{equation}
 \mathcal C_2(\tan\theta)
 =\theta\log(\tan\theta)
   -\int_0^\theta\log(\tan x)\,dx.
 \label{iii:eq:angular-moment}
\end{equation}
The boundary term at zero vanishes because $x\log(\tan x)\to0$.
The tangent difference formula gives
\[
 \tan\frac{\pi}{12}
 =\tan\left(\frac\pi4-\frac\pi6\right)
 =\frac{\sqrt3-1}{\sqrt3+1}=2-\sqrt3.
\]
It remains to evaluate the series in \eqref{iii:eq:log-tangent} at
$\theta=\pi/12$. Let $\chi_{-4}(n)=(-1)^{(n-1)/2}$ for odd $n$.
In the six odd classes modulo $12$, one verifies
\begin{equation}
 2\sin\frac{n\pi}{6}
 =\chi_{-4}(n)+3\,\mathbf1_{3\mid n}\chi_{-4}(n/3).
 \label{iii:eq:character-pell}
\end{equation}
The finite check is
\[
\begin{array}{c|rrrrrr}
 n\bmod12&1&3&5&7&9&11\\ \hline
 2\sin(n\pi/6)&1&2&1&-1&-2&-1\\
 \chi_{-4}(n)&1&-1&1&-1&1&-1\\
 3\mathbf1_{3\mid n}\chi_{-4}(n/3)&0&3&0&0&-3&0
\end{array}
\]
and both sides are periodic modulo $12$ on odd indices.
Absolute convergence permits separation of summands and the substitution $n=3m$
in the second term:
\begin{align*}
 \sum_{\substack{n\geq1\\n\text{ odd}}}
       \frac{\sin(n\pi/6)}{n^2}
 &=\frac12\sum_{n\text{ odd}}\frac{\chi_{-4}(n)}{n^2}
  +\frac32\sum_{m\text{ odd}}\frac{\chi_{-4}(m)}{(3m)^2}\\
 &=\left(\frac12+\frac16\right)G_{\rm Cat}
 =\frac23G_{\rm Cat}.
\end{align*}
Sums without a lower bound in this expression run over positive
integers. Substitution in \eqref{iii:eq:angular-moment}, using
$\log\rho_{\rm P}=-\log\lambda_{\rm P}$, yields
\eqref{iii:eq:catalan-pell}. The first derivative follows from
$\arctan\rho_{\rm P}=\pi/12$; the second from
\[
 \frac{\rho_{\rm P}}{1+\rho_{\rm P}^2}
 =\frac1{\rho_{\rm P}+\rho_{\rm P}^{-1}}=\frac14.
\]
\end{proof}

The coefficient $2/3$ comes from the character decomposition in
\eqref{iii:eq:character-pell}: restriction to multiples of $3$
changes the depth weight by the factor $3^{-2}$. Evaluation
thus brings together the quarter-turn and dodecaphasic subdivision in an
exact moment identity. The coefficients are fixed before any
decimal comparison.

\subsection{Composition of phase and weighted depth}

The depth representation specifies what this evaluation preserves.
On $\ell^2(\mathbb N_{>0})$, define
\[
 N\delta_n=n\delta_n,\qquad U_4\delta_n=i^n\delta_n,
 \qquad Q_4=\frac{U_4-U_4^*}{2i}.
\]
For $0<s<1$, the operator $N^{-2}s^NQ_4$ is trace class, and
\begin{equation}
 \mathcal C_2(s)=\operatorname{Tr}(N^{-2}s^NQ_4).
 \label{iii:eq:pell-depth-trace}
\end{equation}
This equality follows by evaluation on the basis $\delta_n$; the
diagonal entries are $s^n\sin(n\pi/2)/n^2$. The interchange
$U_4\leftrightarrow U_4^*$ reverses $Q_4$ and with it the oriented
moment. Index $n$ and weight $s^n$ remain fixed.

Composition is performed before taking the trace. For
$a\in\mathbb Z/4\mathbb Z$, let $B(s,a)=s^NU_4^a$. Since both
factors are diagonal in the same basis,
\begin{equation}
 B(s,a)B(t,b)=B(st,a+b),\qquad
 B(s,1)^4=s^{4N}.
 \label{iii:eq:pell-depth-composition}
\end{equation}
Phase return preserves the accumulated contraction. In particular,
$B(\rho_{\rm P},1)^4=\rho_{\rm P}^{4N}$, whose eigenvalue at depth
$n$ is $\rho_{\rm P}^{4n}$. This realization gives an operator
meaning to the contractive evaluation; the scalar moment sum
aggregates depths after preserving their index and orientation.
The constitutive arguments $s_\pm=q_\pm^{30}$ remain determined by
the angular pair of \eqref{iii:eq:recover-angular}. The preceding formulas
contain no equality between those arguments and $\rho_{\rm P}$.

\subsection{Dodecaphasic incidence and return functional}

The Barbero--Immirzi contribution uses angular transport and
hexad--octad incidence. Let $\mathsf H\subset\mathsf O$ be a hexad
and an octad with the hexad's marked pair preserved. The relevant flag
sets are
\begin{equation}
 \mathcal F_6=\mathsf H\times\binom{\mathsf H}{2}
 \lhook\joinrel\longrightarrow
 \mathcal F_8=\mathsf O\times\binom{\mathsf H}{2},
 \qquad |\mathcal F_6|=90,\quad |\mathcal F_8|=120.
 \label{iii:eq:barbero-flags}
\end{equation}
The cardinalities are $6\binom62$ and $8\binom62$; the same incidence
preserves the normalized defect
$(90-120)/(24\binom62)=-1/12$. This count identifies the two sectors
used by return, while the calendar determines their
relative multiplicities.

For $m\in\{90,120\}$, the tritic return on a contractive channel is
\begin{equation}
 S_m(q)=\sum_{j\geq0}q^{m+3mj}
       =\frac{q^m}{1-q^{3m}},\qquad 0<q<1.
 \label{iii:eq:barbero-return-series}
\end{equation}
Each term corresponds to the initial height $m$ and $j$ returns of
length $3m$. The geometric series thus preserves both the initial
height and return period. The dodecaphasic calendar
$\mathbb Z/12\mathbb Z$ supplies twelve sectors and a unique oriented
boundary incidence, denoted $\partial_{\rm ret}$. In the free
module of events, its fundamental chain is
\begin{equation}
 c_{12}^{+}
 =\sum_{j\in\mathbb Z/12\mathbb Z}[j,\mathcal F_6]
  -[\partial_{\rm ret},\mathcal F_8].
 \label{iii:eq:barbero-cycle}
\end{equation}
The linear readout $\mathscr L_q$ assigns $S_{90}(q)$ to each generator
$[j,\mathcal F_6]$ and $S_{120}(q)$ to the boundary generator. This yields
\begin{equation}
 \mathscr L_q(c_{12}^{+})=12S_{90}(q)-S_{120}(q).
 \label{iii:eq:barbero-weights}
\end{equation}
The weights $12$ and $-1$ are the chain's oriented multiplicities.
The conjugate orientation satisfies $c_{12}^{-}=-c_{12}^{+}$, with the
same incidence types.

\begin{proposition}[Singular coefficient of the fundamental return]
\label{iii:prop:barbero-pole}
The coefficient of $(1-q)^{-1}$ in
$\mathscr L_q(c_{12}^{+})$ is $1/24$.
\end{proposition}
\begin{proof}
For every integer $m>0$, the factorization
$1-q^{3m}=(1-q)\sum_{j=0}^{3m-1}q^j$ gives
\[
 \lim_{q\uparrow1}(1-q)S_m(q)=\frac1{3m}.
\]
Applying this identity to the functional already determined in
\eqref{iii:eq:barbero-weights} gives
\begin{equation}
 \lim_{q\uparrow1}(1-q)\mathscr L_q(c_{12}^{+})
 =\frac{12}{270}-\frac1{360}=\frac1{24}.
 \label{iii:eq:barbero-pole}
\end{equation}
In the complex coordinate $q-1$, the residue is $-1/24$.
\end{proof}

Singular evaluation verifies a consequence of incidence and tritic
return. The proof order is fixed by
\eqref{iii:eq:barbero-flags}--\eqref{iii:eq:barbero-weights}; the
singular coefficient is computed after the chain has been constructed.

\subsection{Factorization through the constitutive operator}

Retain the lifted angular chart of
Section~\ref{iii:sec:hojas}. To make the angular conversion explicit,
write
\begin{equation}
 A_{\deg}=\frac{180x}{\pi},\qquad
 C^*_{\deg}=\frac{180y}{\pi},\qquad
 q_\pm=e^{-x\mp y},\qquad s_\pm=q_\pm^{30}.
 \label{iii:eq:barbero-angular-lift}
\end{equation}
The orientation considered in \eqref{iii:eq:rchannels} corresponds to
$0<y<x$; sheet interchange is expressed by $y\mapsto-y$ in the
larger domain $x>|y|$. Real exponentials act on these
lifts, which preserve the values of $x\pm y$ during
composition.

The oriented Barbero--Immirzi functional is
\begin{equation}
 \gamma=\frac{\pi A_{\deg}C^*_{\deg}}{180}
   +12\bigl[S_{90}(q_-)-S_{90}(q_+)\bigr]
   -\bigl[S_{120}(q_-)-S_{120}(q_+)\bigr].
 \label{iii:eq:barbero-functional}
\end{equation}
The first component is the angular evaluation; the second is the
difference of the fundamental-chain readout on the conjugate channels.
The arguments are received from the previously constructed angular transport.
The formula determines a dimensionless scalar on that starting
structure.

Let $\psi:(3/4,1)\to(0,1)$ be the inverse of $f$ constructed through
\eqref{iii:eq:cubic}. Define
\begin{equation}
 \mathcal H(s)=\frac{12s^3}{1-s^9}
             -\frac{s^4}{1-s^{12}}
             -\frac{(\log s)^2}{20\pi},\qquad 0<s<1.
 \label{iii:eq:barbero-h}
\end{equation}
This function is real analytic on its domain. The exponents $3,4,9,12$
result from expressing heights $90,120$ and returns $270,360$ in
units of $30$; the logarithmic coefficient is fixed when
the angular contribution is transported.

\begin{theorem}[Oriented constitutive expression of Barbero--Immirzi]
\label{iii:thm:barbero-factorization}
In chart \eqref{iii:eq:barbero-angular-lift},
\begin{equation}
 \gamma=\mathcal H(s_-)-\mathcal H(s_+)
       =\mathcal H\bigl(\psi(r_-)\bigr)
        -\mathcal H\bigl(\psi(r_+)\bigr).
 \label{iii:eq:barbero-factorization}
\end{equation}
On the two-dimensional channel representation,
\begin{equation}
 \gamma=-\operatorname{Tr}
    \left[\mathcal R\,\mathcal H\bigl(\psi(\mathcal V)\bigr)\right],
 \label{iii:eq:barbero-trace}
\end{equation}
where $\operatorname{Tr}$ is the ordinary trace, with
$\operatorname{Tr}I=2$.
\end{theorem}
\begin{proof}
Substitution of $s=q^{30}$ into \eqref{iii:eq:barbero-return-series} gives
\[
 S_{90}(q)=\frac{s^3}{1-s^9},\qquad
 S_{120}(q)=\frac{s^4}{1-s^{12}}.
\]
Moreover, \eqref{iii:eq:barbero-angular-lift} implies
\[
 \log s_\pm=-30(x\pm y)
 =-\frac\pi6(A_{\deg}\pm C^*_{\deg}).
\]
By the difference of squares,
\begin{align}
 \frac{(\log s_+)^2-(\log s_-)^2}{20\pi}
 &=\frac{900\bigl[(x+y)^2-(x-y)^2\bigr]}{20\pi}\\
 &=\frac{180xy}{\pi}
 =\frac{\pi A_{\deg}C^*_{\deg}}{180}.
 \label{iii:eq:barbero-log-term}
\end{align}
Adding these identities proves the first equality in
\eqref{iii:eq:barbero-factorization}. The second uses the inversion
$\psi(f(s_\pm))=s_\pm$ proved in
Section~\ref{iii:sec:constitutiva}.

For the operator expression, the sheet projectors have rank one and
$\mathcal RP_\pm=\pm P_\pm$. Spectral calculus gives
\[
 \mathcal H\bigl(\psi(\mathcal V)\bigr)
 =\mathcal H(s_+)P_++\mathcal H(s_-)P_-.
\]
Multiplying by $\mathcal R$ and taking the ordinary trace gives
$\mathcal H(s_+)-\mathcal H(s_-)$. Its negative is $\gamma$.
\end{proof}

The orientation mark $\mathcal R$ belongs to the functional. If that
mark is held fixed and the channels of $\mathcal V$ are interchanged,
the sign of $\gamma$ changes. Evaluation of the spectrum without its labels preserves
the unordered pair, whereas \eqref{iii:eq:barbero-trace} preserves the
oriented contribution. If the normalized trace
$\tau_2=\operatorname{Tr}/2$ is used, the same formula reads
$\gamma=-2\tau_2[\mathcal R\mathcal H(\psi(\mathcal V))]$.

By contrast, a simultaneous change of representation
$(\mathcal R,\mathcal V)\mapsto
(U\mathcal RU^{-1},U\mathcal VU^{-1})$ preserves the functional:
spectral calculus is transported by conjugation and the trace is
invariant. Interchanging the response relative to a fixed orientation
and jointly transporting response and orientation are therefore
different operations.

This normalization is compatible with the amplifications of
\eqref{iii:eq:logvac}. For
$\mathcal V_m=\mathcal V\otimes I_{9^m}$ and
$\mathcal R_m=\mathcal R\otimes I_{9^m}$,
\begin{equation}
 \gamma=-\frac1{9^m}\operatorname{Tr}
 \left[\mathcal R_m\mathcal H\bigl(\psi(\mathcal V_m)\bigr)\right].
 \label{iii:eq:barbero-amplification}
\end{equation}
Indeed, functional calculus preserves the tensor product and the trace
of $I_{9^m}$ is $9^m$. Representation multiplicity is separated
from the oriented evaluation.

\subsection{Reconstruction from vacuum coordinates}

Identities \eqref{iii:eq:four-identities} express the same
functional through normalized impedance and velocity:
\begin{align}
 \gamma={}&
 \mathcal H\!\left(\psi\!\left(
     \sqrt{\frac{\widehat Z}{\widehat c}}\right)\right)
 -\mathcal H\!\left(\psi\!\left(
     \frac1{\sqrt{\widehat Z\widehat c}}\right)\right).
 \label{iii:eq:barbero-zc}
\end{align}
The positive domain of this reconstruction is described by
\begin{equation}
 1<\widehat Z\widehat c<\frac{16}{9},\qquad
 \frac9{16}<\frac{\widehat Z}{\widehat c}<1.
 \label{iii:eq:barbero-zc-domain}
\end{equation}
Indeed, these inequalities are equivalent to
$3/4<r_\pm<1$ after substitution of the positive roots of
\eqref{iii:eq:four-identities}. For orientation $y>0$, one also retains
$r_-<r_+$, equivalently $\widehat Z<1$. Domain points
are used with the angular compatibility already established by
the TPK; inversion is a subsequent reconstruction of its channels.
The coordinates are the internal ones of \eqref{iii:eq:four}: a change of
units must be undone through \eqref{iii:eq:undo-units} before
$\psi$ is applied.

Interchange $y\mapsto-y$ acts as
$(\widehat Z,\widehat c)\mapsto(\widehat Z^{-1},\widehat c)$ and
produces $\gamma\mapsto-\gamma$. In the balanced case $y=0$,
the channels coincide and \eqref{iii:eq:barbero-factorization} gives
$\gamma=0$. Normalized velocity preserves the product of the
channels; its joint readout with impedance recovers both
responses and the orientation required to evaluate the functional.
In this balanced case $\mathcal V=rI$ has degenerate spectrum:
the mark $\mathcal R$ is retained as representation data, although
the scalar $r$ no longer distinguishes its projectors. For simple spectrum,
by contrast,
$P_+=(\mathcal V-r_-I)/(r_+-r_-)$ and $P_-=I-P_+$.

The central result of this composition is factorization
\eqref{iii:eq:barbero-factorization}. The Catalan moment supplies the
oriented coefficient determining $f$; cubic inversion recovers
each channel's argument; dodecaphasic incidence and the angular pair
determine $\mathcal H$. The Pell scale distinguishes an exact
evaluation of the collection, with the proof of
\eqref{iii:eq:catalan-pell}. Each relation preserves its argument and
operation: the endpoint constant, contractive moment and constitutive
functional are linked evaluations of specified
structures.
